\section{Potências e Raízes}

Potências e raízes são operações frequentes e aparecem em 
fórmulas de desvio padrão e variância. A FX-82MS oferece teclas dedicadas 
para as formas mais comuns dessas operações, e todas assumem que a 
calculadora está no modo \textbf{COMP} \cite{casio2004}.

\subsection{Quadrado e Cubo}

Para elevar um número ao quadrado, utilize a tecla \textbf{$x^2$}, 
localizada no painel central do teclado. Para elevar ao cubo, pressione 
\textbf{SHIFT} seguido de \textbf{$x^2$}, acessando a função secundária 
\textbf{$x^3$}. Ambas as funções são aplicadas ao valor imediatamente 
anterior ao cursor \cite{casio2004}.

Um exemplo de aplicação é o cálculo da variância amostral, que envolve elevar desvios ao quadrado:

\begin{equation}
    s^2 = \frac{(x_1 - \bar{x})^2 + (x_2 - \bar{x})^2 + (x_3 - \bar{x})^2}{n-1}
\end{equation}

\begin{center}
\begin{tabular}{lll}
    \hline
    \textbf{Expressão} & \textbf{Entrada} & \textbf{Resultado} \\
    \hline
    $5^2$ & \texttt{5 $x^2$ =} & \texttt{25} \\
    $3^3$ & \texttt{3 SHIFT $x^2$ =} & \texttt{27} \\
    \hline
\end{tabular}
\end{center}

\subsection{Potência de Expoente Qualquer}

Para calcular uma potência com expoente arbitrário, utilize a tecla
\textbf{$\wedge$}. A sintaxe é: digite a base, pressione a tecla, digite
o expoente e confirme com \textbf{=} \cite{casio2004}.

\begin{center}
\begin{tabular}{lll}
    \hline
    \textbf{Expressão} & \textbf{Entrada} & \textbf{Resultado} \\
    \hline
    $2^{10}$ & \texttt{2 \^{} 10 =} & \texttt{1024} \\
    $3^4$ & \texttt{3 \^{} 4 =} & \texttt{81} \\
    $(-2)^3$ & \texttt{( (-) 2 ) \^{} 3 =} & \texttt{-8} \\
    \hline
\end{tabular}
\end{center}

Atenção: ao calcular potências de bases negativas, utilize parênteses ao
redor da base para garantir que o sinal negativo seja incluído
corretamente na operação. Sem os parênteses, a calculadora pode
interpretar a expressão de forma distinta do esperado \cite{casio2004}.

\subsection{Raiz Quadrada e Cúbica}

A raiz quadrada é acessada diretamente pela tecla 
\textbf{$\sqrt{\phantom{x}}$}, que precede o radicando. A raiz cúbica é 
acessada via \textbf{SHIFT} seguido da tecla de raiz quadrada, ativando 
a função \textbf{$\sqrt[3]{\phantom{x}}$} \cite{casio2004}.

O desvio padrão amostral envolve o cálculo de uma raiz quadrada:

\begin{equation}
    s = \sqrt{\frac{\sum_{i=1}^{n}(x_i - \bar{x})^2}{n-1}}
\end{equation}

\begin{center}
\begin{tabular}{lll}
    \hline
    \textbf{Expressão} & \textbf{Entrada} & \textbf{Resultado} \\
    \hline
    $\sqrt{144}$ & \texttt{$\sqrt{\phantom{x}}$ 144 =} & \texttt{12} \\
    $\sqrt{2}$ & \texttt{$\sqrt{\phantom{x}}$ 2 =} & \texttt{1.414213562} \\
    $\sqrt[3]{27}$ & \texttt{SHIFT $\sqrt{\phantom{x}}$ 27 =} & \texttt{3} \\
    \hline
\end{tabular}
\end{center}

\subsection{Raiz de Índice Qualquer}

Para raízes de índice arbitrário, utilize a função \textbf{$x^{1/y}$}, 
acessada via \textbf{SHIFT} seguido da tecla \textbf{$\wedge$}. A sintaxe 
é: digite o índice da raiz, pressione \textbf{SHIFT} e a tecla, depois 
digite o radicando e confirme com \textbf{=} \cite{casio2004}.

\begin{equation}
    \sqrt[y]{x} = x^{1/y}
\end{equation}

\begin{center}
\begin{tabular}{lll}
    \hline
    \textbf{Expressão} & \textbf{Entrada} & \textbf{Resultado} \\
    \hline
    $\sqrt[4]{625}$ & \texttt{4 SHIFT \^{} 625 =} & \texttt{5} \\
    $\sqrt[5]{32}$ & \texttt{5 SHIFT \^{} 32 =} & \texttt{2} \\
    \hline
\end{tabular}
\end{center}

Ressalta-se que raízes de índice par aplicadas a radicandos negativos 
resultam em erro no domínio real (\texttt{Math ERROR}), pois a FX-82MS 
não opera com números complexos \cite{casio2004}.

\subsection{Combinação Simples}

O cálculo de combinações simples é utilizado na determinação do número de agrupamentos não ordenados de $k$ elementos escolhidos a partir de um total de $n$ elementos. A FX-82MS dispõe da tecla dedicada \textbf{nCr} para realizar esse cálculo diretamente, sem a necessidade de expandir os fatoriais da fórmula \cite{casio2004}:

\begin{equation}
    \binom{n}{k} = C_{n,k} = \frac{n!}{k!(n-k)!}
\end{equation}

A sintaxe de entrada na calculadora é: digite o valor total de elementos ($n$), pressione a tecla \textbf{nCr}, digite a quantidade de elementos agrupados ($k$) e confirme com \textbf{=} \cite{casio2004}.

\begin{center}
\begin{tabular}{lll}
    \hline
    \textbf{Expressão} & \textbf{Entrada} & \textbf{Resultado} \\
    \hline
    $\binom{5}{2}$ & \texttt{5 nCr 2 =} & \texttt{10} \\
    $\binom{10}{3}$ & \texttt{10 nCr 3 =} & \texttt{120} \\
    \hline
\end{tabular}
\end{center}

Essa função constitui o elemento central para a resolução de problemas baseados na \textbf{Distribuição Binomial} no modo \textbf{COMP}, na qual a probabilidade de obter $k$ sucessos em $n$ ensaios independentes é dada por:

\begin{equation}
    P(X = k) = \binom{n}{k} \cdot p^k \cdot (1-p)^{n-k}
\end{equation}

\subsection{Números Negativos e Domínio das Operações}

A FX-82MS opera exclusivamente no domínio dos números reais, o que impõe
restrições a certas combinações de base, expoente e radicando
\cite{casio2004}. Quando uma operação resulta em valor fora desse domínio,
a calculadora exibe \texttt{Math ERROR} e interrompe o cálculo.

As situações mais comuns que produzem esse erro são \cite{casio2004}:

\begin{itemize}
    \item \textbf{Raiz de índice par com radicando negativo} -- por exemplo,
    $\sqrt{-4}$ ou $\sqrt[4]{-16}$, cujos resultados são imaginários;
    \item \textbf{Potência de base negativa com expoente fracionário} --
    por exemplo, $(-8)^{0{,}5}$, equivalente a $\sqrt{-8}$ no reais;
    \item \textbf{Raiz de índice par com expoente fracionário não simplificável}
    -- casos em que a calculadora não consegue determinar um resultado real.
\end{itemize}

Raízes de índice \emph{ímpar} aplicadas a radicandos negativos, por outro
lado, são válidas no domínio real e calculadas corretamente:

\begin{center}
\begin{tabular}{lll}
    \hline
    \textbf{Expressão} & \textbf{Entrada} & \textbf{Resultado} \\
    \hline
    $\sqrt[3]{-8}$  & \texttt{SHIFT $\sqrt{\phantom{x}}$ (-) 8 =} & \texttt{-2} \\
    $(-2)^3$ & \texttt{( (-) 2 ) \^{} 3 =} & \texttt{-8} \\
    $(-2)^4$ & \texttt{( (-) 2 ) \^{} 4 =} & \texttt{16} \\
    $\sqrt{-4}$ & \texttt{$\sqrt{\phantom{x}}$ (-) 4 =} & \texttt{Math ERROR} \\
    \hline
\end{tabular}
\end{center}

Ao encontrar \texttt{Math ERROR}, pressione \textbf{AC} para limpar o
visor e revisar a expressão antes de recalcular.

\subsection{Exemplos Práticos}

Esta subseção apresenta aplicações de potências e raízes com
exemplos contextualizados em Probabilidade e Estatística.

\subsubsection*{Exemplo 1 — Desvio padrão amostral}

Dadas as notas $6$, $8$ e $10$, com média $\bar{x} = 8$, o desvio padrão
amostral é:

\begin{equation}
    s = \sqrt{\frac{(6-8)^2 + (8-8)^2 + (10-8)^2}{3-1}}
      = \sqrt{\frac{4 + 0 + 4}{2}} = \sqrt{4} = 2
\end{equation}

O cálculo é realizado em etapas, aproveitando a memória \textbf{Ans}
para encadear os resultados intermediários:

\begin{center}
\begin{tabular}{lll}
    \hline
    \textbf{Etapa} & \textbf{Entrada} & \textbf{Resultado} \\
    \hline
    Soma dos quadrados & \texttt{( 6 - 8 ) $x^2$ + ( 8 - 8 ) $x^2$ + ( 10 - 8 ) $x^2$ =} & \texttt{8} \\
    Variância amostral & \texttt{Ans / ( 3 - 1 ) =} & \texttt{4} \\
    Desvio padrão & \texttt{$\sqrt{\phantom{x}}$ Ans =} & \texttt{2} \\
    \hline
\end{tabular}
\end{center}

\subsubsection*{Exemplo 2 — Variância populacional}

Uma pesquisa registrou os seguintes tempos de atendimento em minutos:
$2$, $4$, $4$, $4$, $5$, $5$, $7$ e $9$. A média é $\bar{x} = 5$ e
a variância populacional é:

\begin{equation}
    s^2 = \frac{\sum_{i=1}^{n}(x_i - \bar{x})^2}{n}
             = \frac{9 + 1 + 1 + 1 + 0 + 0 + 4 + 16}{8}
             = \frac{32}{8} = 4
\end{equation}

\begin{center}
\begin{tabular}{lll}
    \hline
    \textbf{Etapa} & \textbf{Entrada} & \textbf{Resultado} \\
    \hline
    Soma dos quadrados & \texttt{( 2 - 5 ) $x^2$ + ( 4 - 5 ) $x^2$ +} & \\
                       & \texttt{( 4 - 5 ) $x^2$ + ( 4 - 5 ) $x^2$ +} & \\
                       & \texttt{( 5 - 5 ) $x^2$ + ( 5 - 5 ) $x^2$ +} & \\
                       & \texttt{( 7 - 5 ) $x^2$ + ( 9 - 5 ) $x^2$ =} & \texttt{32} \\
    Variância & \texttt{Ans / 8 =} & \texttt{4} \\
    Desvio padrão & \texttt{$\sqrt{\phantom{x}}$ Ans =} & \texttt{2} \\
    \hline
\end{tabular}
\end{center}

\subsubsection*{Exemplo 3 — Coeficiente de variação}

O coeficiente de variação ($CV$) expressa o desvio padrão como percentual
da média, permitindo comparar a dispersão de conjuntos com unidades ou
magnitudes diferentes:

\begin{equation}
    CV = \frac{s}{\bar{x}} \times 100
\end{equation}

Considere duas turmas com os seguintes resultados em uma avaliação:

\begin{itemize}
    \item Turma A: $\bar{x} = 8{,}0$ e $s = 2{,}0$
    \item Turma B: $\bar{x} = 6{,}0$ e $s = 1{,}8$
\end{itemize}

\begin{center}
\begin{tabular}{lll}
    \hline
    \textbf{Etapa} & \textbf{Entrada} & \textbf{Resultado} \\
    \hline
    $CV$ Turma A & \texttt{2 / 8 * 100 =}   & \texttt{25} \\
    $CV$ Turma B & \texttt{1.8 / 6 * 100 =} & \texttt{30} \\
    \hline
\end{tabular}
\end{center}
